\documentstyle[%


righttag%


,11pt%


,amssymb%


,russian%               remove in Eng.


,izvru%                 Set izven% in Eng.


% ,izvru-ks             for brief communications


]{amsart}





%       Math definitions


\newcommand{\eqdef}{\overset{\mathrm{def}}{=}{}}


\newcommand{\spc}[1]{\bold #1}


\newcommand{\vraisup}{\operatornamewithlimits{vrai\,sup}}


\newcommand{\tts}[1]{}


\renewcommand{\baselinestretch}{0.96}





%\hoffset=-15mm%                  For Eng.


\setcounter{page}{28}


\god{2000}


\nomer{4~(455)} %                        Remove in Eng.


%\nomer{Vol.\,44, No.\,4}%               Set in Eng.


%\str{\pageref{begin}--\pageref{end}}%  Set in Eng.


\udk{517.929\,:\,517.925}





\author{\it Ж.\,С.\,П.\,Мунембе}


% NB:   ^^ remove in Eng.


\title{К вопросу об асимптотическом поведении


решений системы линейных функционально-дифференциальных уравнений с


запаздыванием и периодическими параметрами}


\thanks{Работа выполнена при финансовой поддержке


Capacity Building Project (Credit 2436, Eduardo Mondlane University


--- Mozambique).}





\date{}


%\pagestyle{myheadings}%         Set in Eng.


\pagestyle{plain} %              Remove in Eng.


\begin{document}


%\thispagestyle{plain}%          Set in Eng.


\maketitle


\label{begin}





\section{Постановка задачи}





Будем предполагать в задаче Коши для системы функционально-дифференциальных


уравнений с запаздывающим аргументом


\begin{gather}


\Dot{x}(t)-P(t)x[h(t)]=f(t), \quad t\in [0,\infty), \quad


x(\xi)=\varphi(\xi), \ \; \text{если} \ \; \xi <0, \\


%


x(0)={\alpha},


\end{gather}


что столбцы $n\times n$-матрицы $P:[0,{\infty}){\to}{{\bold R}^{n\times


n}}$ и функция $f:[0,{\infty}){\to}{{\bold R}^n}$ периодические с


периодом $T>0$ и суммируемы на периоде; запаздывание


$h:[0,{\infty}){\to}R$ имеет вид $h(t)=t-{\Delta}(t)$,


$0{\le}{\Delta}(t){\le}T$, где функция


${\Delta}:[0,{\infty}){\to}[0,T]$ измерима и периодична с периодом


$T$; начальная функция ${\varphi}:[-T,0]{\to}{{\bold


R}^n}$ такова, что функция ${\varphi}^h:[0,{\infty}){\to}{{\bold


R}^n}$, определенная равенством


$$


{\varphi}^h(t)=


\begin{cases}


0, & \text{если } h(t){\ge}0;\\


{\varphi}[h(t)], & \text{если } h(t)<0,


\end{cases}


$$


измерима и ограничена


в существенном на $[0,{\infty})$.





Всюду ниже через ${\spc L}={\spc L}[0,b]$ будем обозначать


пространство суммируемых функций $f:[0,b]\to {{\bold R}^n}$, а через


${\spc C}={\spc C}[0,b]$ --- пространство непрерывных функций


$x:[0,b]\to {{\bold R}^n}$,


$\|x\|_{\bold C}=\max


\limits_{t\in[0,b]} \|x(t)\|_{{\bold R}^n}$.





Используя стандартное обозначение \cite{1}


$$


x_h(t)=


\begin{cases}


x[h(t)], & \text{если } h(t){\ge}0;\\


0, & \text{если } h(t)<0,


\end{cases}


$$


запишем систему (1) в виде


\begin{equation}


({\cal L}x)(t)\equiv{\Dot x}(t)-P(t)x_h(t)=f(t)+P(t){\varphi}^h(t),


\quad


t\in [0,{\infty}).


\end{equation}


Как известно (\cite{1}, \cite{2}), решение $x(t,{\alpha})$ задачи


(1),\;(2) ((3),\;(2)) представимо в виде


$$


x(t)=C(t,0){\alpha}+\int\limits_0^tC(t,s)P(s){\varphi}^h(s)ds+


\int\limits_0^tC(t,s)f(s)ds,\quad t\in [0,{\infty }),


$$


где $C(t,s)$ --- матрица Коши оператора ${\cal L}$.


\medskip





{\bf Определение.}


Будем говорить, что решение $x(t,{\alpha})$ задачи (1), (2)


стабилизируется к периодической функции $y:[0,{\infty})\to {{\bold


R}^n}$ периода $T$, если


$$


\lim_{{\nu}\to \infty


}\,\max_{t\in [({\nu}-1)T,{\nu}T]}\| x(t,{\alpha})-y(t)


\|_{{\bold R}^n}=0.


$$





В предлагаемой статье будут установлены условия стабилизируемости решения


задачи (1), (2) к периодической функции и даны оценки скорости


стабилизируемости.


Для случая $h:[0,T]\to [0,T]$ такие условия можно найти в \cite{3}.





\section{Свойства матрицы Коши}





{\bf Утверждение.} {\it Матрица Коши операции ${\cal L}$ обладает


свойством периодичности}


$$


C(t+T,s+T)=C(t,s),\quad


0{\le}s{\le}t<{\infty}.


$$





{\bf Доказательство}. Функция $h$ обладает


свойством $h(t+T)=h(t)+T$. Обозначим через ${\chi}_h(t,s)$


характеристическую функцию множества $\{(t,s)\in [0,\infty )\times


[0,\infty ):0\le s\le h(t)< \infty \}$. Так как справедливо


равенство $x_h(t)={\chi_h}(t,0)x(0)+\int\limits_0^t{\chi}_h(t,s){\Dot


x}(s)ds$, то, подставив правую часть этого представления в (3),


получим


$$


{\Dot x}(t)-\int\limits_0^tP(t){\chi}_h(t,s){\Dot


x}(s)ds-P(t){\chi}_h(t,0)x(0)= f(t)+P(t){\varphi}^h(t).


$$


Обозначим


$K(t,s)\eqdef P(t){\chi}_h(t,s)$,


$f_1(t)\eqdef f(t)+P(t){\varphi}^h(t)+P(t){\chi}_h(t,0)x(0)$ и


$z(t)={\Dot x}(t)$.


Тогда предыдущее уравнение примет вид


$$


z(t)-\int\limits_0^tK(t,s)z(s)ds=f_1(t),\quad t\in [0,{\infty}).


$$


Равенство


$$


K(t+T,s+T)=K(t,s),\quad 0\le s\le t<{\infty},


$$


проверяется непосредственно.





Обозначим через $K$ оператор, действующий из ${\spc L}[0,b]$ в ${\spc


L}[0,b]$ по правилу


$$


(Kz)(t)=\int\limits_0^tK(t,s)z(s)ds.


$$


Определим итерированные ядра. Считая $K_1(t,s)\equiv K(t,s)$, для


$m$-й степени $K$ имеем $(K^my)(t)=\int\limits_0^tK_m(t,s)y(s)ds$,


где $K_m(t,s)$ --- $m$-е итерированное ядро:


$K_m(t,s)=\int\limits_s^tK_{m-1}(t,\tau)K(\tau,s)d\tau$,


$m=2,3,\dotsc$\; Покажем, что каждое итерированное ядро периодично


по совокупности переменных.





Действительно, пусть $m=2$. Тогда


$$


K_2(t,s)=\int\limits_s^tK(t,\tau)K(\tau ,s)d\tau


$$


и


\begin{gather*}


K_2(t+T,s+T)=\int\limits_{s+T}^{t+T}P(t+T){\chi}_h


(t+T,{\tau})P({\tau}){\chi}_h


({\tau},s+T)d{\tau}=\\


%


=\int\limits_s^tP(t){\chi}_h(t{+}T,\tau{+}T)P(\tau{+}T)


{\chi}_h(\tau{+}T,s{+}T)


d\tau=\int\limits_s^tP(t){\chi}_h(t,\tau)P(\tau){\chi}_h(\tau,s)d{\tau}=


K_2(t,s).


\end{gather*}





Далее, по индукции легко проверить, что для любого $m>2$ имеет место


равенство


$$


K_m(t+T,s+T)=K_m(t,s),\quad 0\le s\le t<{\infty}.


$$





В ряде


\begin{equation}


\sum_{m=1}^{\infty}K_m(t,s),\quad 0\le s\le t<{\infty},


\end{equation}


итерированные ядра $K_m$ будем считать элементами пространства


${\spc L}_{1,{\infty}}^b$ измеримых функций $K:[0,b]{\times}[0,b]\to 


{\bold R}^{n{\times}n}$ таких, что


$$


\| K(\cdot,\cdot)\|_{\bold L_{1,{\infty}}^b}


\eqdef


\| K(\cdot,\cdot)\|_{1,{\infty}}^b \eqdef


\vraisup_{s\in [0,b]}{\int\limits_s^b\|K(t,s)\|


dt}<{\infty},$$


где $\|\cdot\|$ --- норма матрицы,


согласованная с $\|\cdot\|_{{\bold R}^n}$.





Обозначим через $p(t)$ такую $T$-периодическую суммируемую на периоде функцию,


что $p(t)\ge \|P(t)\|$, $t\in [0,T]$.





Построим для этого ряда мажоранту. Тогда


$$


\|K_m(t,s)\|\le p(t)\frac


{\Big[\int\limits_0^bp(\tau )d{\tau}\Big]^{m-1}}


{(m-1)!}, \quad m=1,2,\dotsc,


$$


и


$$


\|K_m(\cdot,\cdot)\|_{1,{\infty}}^b\le\frac


{\Big[\int\limits_0^bp(\tau )d{\tau}\Big]^m}{(m-1)!},


$$


что гарантирует сходимость


ряда (4) в пространстве ${\spc L}_{1,{\infty}}^b$.





Обозначая через $R(t,s)$ сумму ряда (4), имеем в силу произвольности $b$


$$


R(t,s)=R(t+T,s+T),\quad 0\le s\le t<{\infty}.


$$


Как известно, $R(t,s)$ является


резольвентой ядра $K(t,s)$: при любой правой части $f$ решение уравнения


$z-Kz=f$ представимо в виде $z=f+Rf$, где $R$ --- оператор,


действующий из ${\spc L}[0,b]$ в ${\spc L}[0,b]$ по правилу


$$


(Rf)(t)=\int\limits_0^tR(t,s)f(s)ds.


$$





Далее, из известных \cite{4} соотношений


\begin{gather*}


R(t,s)=C_t'(t,s), \quad 0\le s\le t<\infty ,\\


%


C(t,s)=E+\int\limits_s^tC_{\tau}'(\tau ,s)d\tau ,


\quad 0\le s\le t<\infty ,


\end{gather*}


выражающих связь между резольвентным ядром $R(t,s)$ для ядра


$K(t,s)$ и матрицей Коши $C(t,s)$, имеем


$$


R(t+T,s+T)=C_t'(t+T,s+T).


$$


Отсюда следует


$$


C(t{+}T,s{+}T)=E+\int\limits_{s+T}^{t+T}C_{\tau}'({\tau},s{+}T)d{\tau}=


E+\int\limits_s^tC_{\tau}'({\tau}{+}T,s{+}T)d{\tau}=


E+ \int\limits_s^tC_{\tau}'({\tau},s)d{\tau}


=C(t,s).\ \square


$$





\section{Построение вспомогательной задачи}





Для нахождения условий стабилизируемости построим некоторую вспомогательную


задачу.


Сначала рассмотрим последовательное построение


решения задачи (3),\;(2) ((1),\;(2)).





Разобьем интервал $[0,{\infty})$ точками


$0<T<\dotsb<{\nu}T<({\nu}+1)T<\dotsb$ Обозначим через


$C_{\nu}(t,s)=\{ C_{ij}^{\nu}(t,s)\}_{i,j=1}^n$, $\nu =1,2,\dotsc$,


сужение матрицы Коши $C(t,s)$ на множество


$$


{\cal D_{\nu}}=\{(t,s):t\in [(\nu -1)T,\nu T], \ s\in [0,t]\},


$$


а через


$x_{\nu}(t)$ --- сужение решения задачи (1), (2), определенное на


промежутке $[({\nu}-1)T,{\nu}T]$, ${\nu}=1,2,\dotsc$\; Через


${\chi}_{\nu}(t)$ обозначим характеристическую функцию отрезка


$[({\nu}-1)T,{\nu}T]$, ${\nu}=1,2,\,\dotsc$\; В этих обозначениях


решение задачи (1), (2) может быть представлено в виде


$$


x(t,{\alpha})=\sum_{\nu


=1}^{(N+1)T}x_{\nu}(t){\chi}_{\nu}(t), \quad N=1,2,\,\dotsc


$$


Функция


$x_{\nu}(t)$, $t\in [({\nu}-1)T,{\nu}T]$, ${\nu}=1,2,\dotsc$, является


решением задачи Коши


\begin{gather}


({\cal L_{\nu}}x_{\nu})(t){\equiv}{\Dot


x_{\nu}}(t)-


P(t)[({\Upsilon}_{{\nu}-1}^hx_{\nu})+({\Gamma}_{{\nu}-1}^hx_{{\nu}-1})](t)


=


f(t), \quad t\in [(\nu -1)T,\nu T], \tag{$5_\nu$} \\


%


x_{\nu}(({\nu}-1)T)={\alpha}_{\nu}.\tag{$6_\nu$}


\end{gather}


Здесь ${\alpha}_1={\alpha}$,


${\alpha}_{\nu}=x_{{\nu}-1}(({\nu}-1)T)$, ${\nu}=2,3,\dotsc$,


$x_0(t){\equiv}{\varphi}(t)$,


$({\Gamma}_0^hx_0)(t)={\varphi}^h(t)$,


$$


({\Upsilon}_{\nu


-1}^hx_{\nu})(t)=


\begin{cases}


x_{\nu}[h(t)], &


\text{если } h(t) \in [(\nu -1)T,{\nu}T];\\


0, & \text{если } h(t)\in [(\nu -2)T,({{\nu}-1})T],


\end{cases}


\quad t\in [({\nu}-1)T,{\nu}T],


$$


и


$$


({\Gamma}_{\nu -1}^hx_{{\nu}-1})(t)=


\begin{cases}


0, & \text{если } h(t) \in [(\nu -1)T,{\nu}T];\\


x_{{\nu}-1}[h(t)], & \text{если } h(t)\in [(\nu


-2)T,({{\nu}-1})T],


\end{cases}


\quad t\in [({\nu}-1)T,{\nu}T].


$$





Сужение $x_1(t)$ решения задачи (3), (2) имеет вид


$$


x_1(t)=C_1(t,0)x_1(0)+\int\limits_0^tC_1(t,s)P(s){\varphi}^h(s)ds


+\int\limits_0^tC_1(t,s)f(s)ds,\quad t\in [0,T].


$$


В силу известной формулы Коши (\cite{5}, с.\,92--95)


на втором промежутке $[T,2T]$ сужение решения задачи (3), (2) представимо в


виде


$$


x_2(t)=C_2(t,T)x_1(T)+\int\limits_T^tC_2(t,s)P(s)({\Gamma}_1^hx_1)(s)ds+


\int\limits_T^tC_2(t,s)f(s)ds,


\quad t\in [T,2T].


$$


Аналогично для сужения $x_{\nu}(t)$ решения задачи (3), (2)


получаем представление


\begin{multline*}


x_{\nu}(t)=C_{\nu}(t,(\nu -1)T)x_{\nu -1}((\nu -1)T)+


\int\limits_{(\nu -1)T}^tC_{\nu}(t,s)P(s)({\Gamma}_{{\nu}-1}^hx_{\nu


-1})(s)ds+ \\


%


+\int\limits_{(\nu -1)T}^tC_{\nu}(t,s)f(s)ds,\quad


t\in [(\nu -1)T,\nu T].


\end{multline*}





Упомянутая выше вспомогательная задача строится на основе последовательности


задач ($5_{\nu}$)--($6_{\nu}$) следующим образом.





Построим на отрезке $[0,T]$ последовательность функций


\begin{alignat*}{2}


{\ov x_1}(t)\eqdef & x_1(t), &\quad t &\in [0,T], \\


%


{\ov x_2}(t)\eqdef & x_2(t+T), &\quad t &\in [0,T], \\


%


&\dotsc\dotsc\dotsc & \quad &\ \\


%


{\ov x_{\nu}}(t)\eqdef & x_{\nu}(t+(\nu -1)T), &\quad t &\in [0,T],


\end{alignat*}


где $x_{\nu}(\cdot)$ --- решение задачи ($5_{\nu}$)--($6_{\nu}$).


Определим оператор ${\Gamma}^{h+T}:{\spc C}[0,T]\to {\spc L}[0,T]$


равенством


$$


({\Gamma}^{h+T}z)(t)=


\begin{cases}


0, & \text{если } h(t)+T\notin [0,T];\\


z[h(t)+T], & \text{если } h(t)+T\in [0,T],


\end{cases}


\quad t\in [0,T].


$$


Покажем, что\addtocounter{equation}{2}


\begin{gather}


{\ov x_1}(t)=C(t,0){\alpha}+\int\limits_0^tC(t,s)P(s)


{\varphi}^h(s)ds+


\int\limits_0^tC(t,s)f(s)ds, \\


%


{\ov x_{\nu}}(t)=C(t,0){\ov x_{\nu -1}}(T)+


\int\limits_0^tC(t,s)P(s)({\Gamma}^{h+T}{\ov x_{\nu -1}})(s)ds+


\int\limits_0^tC(t,s)f(s)ds,\quad


{\nu}=2,3,\dotsc \tag{$8_\nu$}


\end{gather}





При $\nu =2$ равенство ($8_\nu$) справедливо, т.\,к.


в силу периодичности матрицы Коши имеем


\begin{gather*}


C_{\nu}(t,s)=C_{\nu -1}(t+T,s+T)=C(t,s), \quad \nu


=1,2,\dotsc,\quad C_0(t,s)\eqdef C(t,s), \qquad


(t,s)\in {\cal D_{\nu}},\\


%


{\ov x_2}(t)\eqdef x_2(t+T)=


C_2(t+T,T)x_2(T)+\int\limits_T^{t+T}C_2(t+T,s)P(s)({\Gamma}_1^hx_1)(s)ds+


\int\limits_T^{t+T}C_2(t+T,s)f(s)ds= \\


%


=C(t,0){\ov x_1}(T)+


\int\limits_0^tC_2(t+T,s+T)P(s+T)({\Gamma}_1^hx_1)(s+T)ds+


\int\limits_0^tC_2(t+T,s+T)f(s+T)ds= \\


%


=C(t,0){\ov x_1}(T)+


\int\limits_0^tC(t,s)P(s)({\Gamma}^{h+T}{\ov x_1})(s)ds+


\int\limits_0^tC(t,s)f(s)ds.


\end{gather*}


Пусть ($8_\nu$) выполнено при некотором ${\nu}$. Убедимся, что


имеет место ($8_{\nu+1}$). Действительно,


\begin{gather*}


{\ov x_{{\nu}+1}}(t)\eqdef x_{{\nu}+1}(t+{\nu}T)=\\


%


=C_{{\nu}+1}(t+{\nu}T,{\nu}T)x_{{\nu}}({\nu}T)+


\int\limits_{{\nu}T}^


{t+{\nu}T}C_{{\nu}+1}(t+{\nu}T,s)


P(s)({\Gamma}_{{\nu}}^hx_{{\nu}})(s)ds+


\int\limits_{{\nu}T}^{t+{\nu}T}C_{{\nu}+1}(t+{\nu}T,s)f(s)ds=\\


%


=C(t,0){\ov x_{{\nu}}}(T)+


\int\limits_0^tC_{{\nu}+1}(t+{\nu}T,s+{\nu}T)


P(s+{\nu}T)({\Gamma}_{{\nu}}^hx_{{\nu}})(s+{\nu}T)ds+\\


%


+\int\limits_0^tC_{{\nu}+1}(t+{\nu}T,s+{\nu}T)f(s+{\nu}T)ds=\\


%


=C(t,0){\ov x_{{\nu}}}(T)+


\int\limits_0^tC(t,s)P(s)({\Gamma}^{h+T}{\ov x_{{\nu}}})(s)ds+


\int\limits_0^tC(t,s)f(s)ds.


\end{gather*}





Обозначив


$$


g(t)=\int\limits_0^tC(t,s)f(s)ds,


$$


запишем равенство ($8_{\nu}$) в виде\addtocounter{equation}{1}


\begin{equation}


{\ov x_{\nu}}(t)=C(t,0){\ov x_{\nu -1}}(T)+


\int\limits_0^tC(t,s)P(s)({\Gamma}^{h+T}{\ov x_{\nu -1}})(s)ds+g(t),


\quad {\nu}=2,3,\dotsc


\end{equation}


Здесь функция ${\ov x_1}(t)=x(t,{\alpha})$, $t\in [0,T]$, определена


соотношением (7).





Соотношение (9) определяет процесс построения простых итераций для уравнения


\begin{equation}


z(t)=C(t,0)z(T)+


\int\limits_0^tC(t,s)P(s)({\Gamma}^{h+T}z)(s)ds+g(t),


\quad t\in [0,T].


\end{equation}





Заметим, что уравнение (10) относится к так называемым нагруженным


интегральным


уравнениям с отклоняющимся аргументом. Заметим также, что любое решение


уравнения (10) удовлетворяет периодическому условию $z(0)=z(T)$.





Задача о стабилизируемости решения задачи (3), (2) к периодической функции


свoдится к задаче о сходимости


итерационного процесса для вспомогательного уравнения (10)


(см. п.5\,).


В случае $h:[0,T]\to [0,T]$ уравнение (10) принимает вид


$$


z(t)=C(t,0)z(T)+g(t)


$$


и, как легко видеть, итерационный процесс (9) сходится,


если спектральный радиус оператора монодромии $C(T,0):{\bold


R}^n\to {\bold R}^n$ меньше единицы.





\section{Условие сходимости итераций для вспомогательного


уравнения}





Для формулировки условий сходимости итерационого процесса (9) запишем


(10) в виде


\begin{equation}


z=(A+B)z+g,


\end{equation}


где линейные операторы $A,B:{\spc C}[0,T]\to {\spc C}[0,T]$ определяются


равенствами


\begin{gather*}


(Az)(t)=C(t,0)z(T), \\


%


(Bz)(t)=\int\limits_0^tC(t,s)P(s)({\Gamma}^{h+T}z)(s)ds.


\end{gather*}





Обозначим через ${\rho}={\rho}(A+B)$ спектральный радиус оператора


$$


(A+B):{\spc C}[0,T]\to {\spc C}[0,T].


$$


Как известно, при выполнении условия


\begin{equation}


{\rho}<1


\end{equation}


уравнение (10) имеет единственное решение $z^*\in {\spc C}[0,T]$ и


итерационный процесс (9) сходится в смысле нормы


пространства ${\spc C}[0,T]$ к $z^*$, т.\,е. $\|{\ov


x_{\nu}}-z^* \|_{\bold C[0,T]}\to 0$.





Явную оценку скорости такой сходимости дает





\begin{thm}


Пусть выполнено условие $(12)$. Тогда


для каждого ${\epsilon}>0$ такого, что ${\delta}


\eqdef {\rho}+{\epsilon}<1$, найдется такое натуральное


$k_0\eqdef k_0({\delta})$, что имеет место оценка


$$


\|{\ov x_{\nu}}-z^*\|_{\bold C[0,T]}


\le \frac{{\delta}^{{\nu}-1}}{1-{\delta}} \frac{M}{m}


\|{\ov x_1}-(A+B){\ov x_1}-g\|_{\bold C[0,T]},


$$


где $m\eqdef {\delta}^{k_0-1}$,


$M\eqdef {\delta}^{k_0-1}+{\delta}^{k_0-2}\|A+B\|+\dotsb+


\|(A+B)^{k_0-1}\|$.


\end{thm}





{\bf Доказательство}. Пусть ${\rho}={\rho}(A+B)=


\lim\limits_{k\to \infty }\sqrt[\uproot{2}k]{\| (A+B)^k\|}$


и ${\rho}<1$. Задаем ${\epsilon}>0$ такое, что ${\delta}


\eqdef {\rho}+{\epsilon}<1$.


Используя


конструкцию, описанную в (\cite{6}, с.\,15), введем


в пространство ${\spc C}[0,T]$ такую эквивалентную


норму $\|\cdot\|_*$, при которой норма линейного оператора


$A+B$ будет отличаться от спектрального радиуса не более чем на ${\epsilon}$.


А именно, в силу определения спектрального радиуса найдется такое


$k_0=k_0({\delta})$, что


$$


\sqrt[\uproot{3}{k_0}]{\| (A+B)^{k_0}\|}\le {\delta}.


$$


В этом доказательстве для простоты записи будем опускать индекс


${\spc C}[0,T]$ у нормы $\|\cdot\|_{\bold C[0,T]}$.





Положим


$$


\|{\ov x}\|_*={\delta}^{k_0-1}\|{\ov x}\|+


{\delta}^{k_0-2}\|(A+B){\ov


x}\|+\,\dotsc\,+\| (A+B)^{k_0-1}{\ov x}\|.


$$


Очевидно


$$


m\|\ov x\|\le \|{\ov x}\|_*\le M\|{\ov x}\|,


$$


т.\,е. нормы $\|\cdot\|$ и $\|\cdot\|_*$


эквивалентны. Покажем, что 


$\|A+B\|^*\le {\delta}$ в новой норме (здесь


$\|A+B\|^*\eqdef \sup\limits_{\|x\|_*\le 1}\|(A+B)x\|_*$ ).


Действительно,


\begin{gather*}


\|(A+B){\ov x}\|_*={\delta}^{k_0-1}\|(A+B){\ov x}


\|+{\delta}^{k_0-2}\|(A+B)^2{\ov


x}\|+\dotsb+ \|(A+B)^{k_0}{\ov x}\|\le \\


%


\le {\delta}^{k_0-1}\|(A+B){\ov


x}\|+ {\delta}^{k_0-2}\|(A+B)^2{\ov


x}\|+\dotsb+ {\delta}\|(A+B)^{k_0-1}{\ov


x}\|+ \|(A+B)^{k_0}\|\|{\ov x}\|\le \\


%


\le {\delta}^{k_0-1}\|(A+B){\ov


x}\|+ {\delta}^{k_0-2}\|(A+B)^2{\ov


x}\|+\dotsb+ {\delta}\|(A+B)^{k_0-1}{\ov


x}\|+ {\delta}^{k_0}\|{\ov x}\|= \\


%


={\delta}[{\delta}^{k_0-1}\|{\ov x}\|+


{\delta}^{k_0-2}\|(A+B){\ov x}\|+\dotsb+


\|(A+B)^{k_0-1}{\ov x}\|]={\delta}\|{\ov x}\|_*,


\end{gather*}


т.\,е.


$$


\|A+B\|_*\le {\delta}.


$$


Таким образом, оператор $A+B$ является оператором сжатия.


В таком случае при любом начальном элементе ${\ov x_1}(\cdot,{\alpha})$


последовательные приближения (9) сходятся к элементу


$z^*\in {\spc C}[0,T]$, при этом


\begin{equation}


\|{\ov x_{\nu}}-z^*\|_*\le 


{\delta}^{{\nu}-1}


\|{\ov x_1}-z^*\|_*


\end{equation}


и


\begin{equation}


\|{\ov x_1}-z^*\|_*\le \frac{1}{1-{\delta}}


\|{\ov x_1}-(A+B){\ov x_1}-g\|_*.


\end{equation}


В силу (13) и (14) имеем


$$


\|{\ov x_{\nu}}-z^*\|_*\le 


\frac{{\delta}^{{\nu}-1}}{1-{\delta}}\|{\ov x_1}-(A+B){\ov x_1}-g


\|_*.


$$


Используя эквивалентность норм $\|\cdot\|$ и


$\|\cdot\|_*$, получаем


\begin{gather*}


\|{\ov x_{\nu}}-z^*\|\le \frac{1}{m}


\|{\ov x_{\nu}}-z^*\|_*\le \frac{1}{m}


\frac{{\delta}^{{\nu}-1}}{1-{\delta}}\|{\ov x_1}-(A+B){\ov x_1}-g\|_*= \\


%


=\frac{{\delta}^{{\nu}-1}}{1-{\delta}}\frac{1}{m}


({\delta}^{k_0-1}\|{\ov x_1}-(A+B){\ov x_1}-g\|+


{\delta}^{k_0-2}\|(A+B)({\ov x_1}-(A+B){\ov x_1}-g)\|+\dotsb+\\


%


+\|(A+B)^{k_0-1}({\ov x_1}-(A+B){\ov x_1}-g)\|)\le \\


%


\le \frac{{\delta}^{{\nu}-1}}{1-{\delta}}\frac{1}{m}


({\delta}^{k_0-1}+{\delta}^{k_0-2}


\|A+B\|+\dotsb+\|(A+B)^{k_0-1}


\|)\|{\ov x_1}-(A+B){\ov x_1}-g\|=\\


%


=\frac{{\delta}^{{\nu}-1}}{1-{\delta}}


\frac{M}{m}\|{\ov x_1}-(A+B){\ov x_1}-g\|.\qquad\square


\end{gather*}





\section{Условие и оценка скорости стабилизируемости}





\begin{thm}


Пусть спектральный радиус ${\rho}$ оператора $A+B$ меньше


единицы и $y:[0,{\infty})\to {\bold R}^n$ --- $T$-периодическое


продолжение решения $z^*\in \bold C[0,T]$ уравнения $(11)$.


Тогда для любого такого ${\varepsilon}>0$, что ${\delta}


\eqdef {\rho}(A+B)+{\varepsilon}<1$, найдется


натуральное $k_0=k_0({\delta})$, при котором имеет место оценка


\begin{equation}


\max_{t\in [({\nu}-1)T,{\nu}T]}\| x(t,{\alpha})-y(t)


\|_{{\bold R}^n} \le 


\frac{{\delta}^{{\nu}-1}}{1-{\delta}}


\frac{M}{m}\|x(\cdot,{\alpha})-(A+B)x(\cdot,{\alpha})-g


\|_{\bold C[0,T]}.


\end{equation}


\end{thm}





\begin{pf} Действительно,


\begin{multline*}


\max_{t\in [({\nu}-1)T,{\nu}T]}\| x(t,{\alpha})-y(t)


\|_{{\bold R}^n}=


\max_{t\in [({\nu}-1)T,{\nu}T]}\| x_{\nu}(t)-z^*(t)


\|_{{\bold R}^n}= \\


%


=\max_{t\in [0,T]}\| x_{\nu}(t+({\nu}-1)T)-z^*(t)


\|_{{\bold R}^n}=


\max_{t\in [0,T]}\| {\ov x_{\nu}}(t)-z^*(t)


\|_{{\bold R}^n}.


\end{multline*}


Отсюда в силу теоремы 1 имеем


$$


\max_{t\in [0,T]}\| {\ov x_{\nu}}(t)-z^*(t)


\|_{{\bold R}^n}\le 


\frac{{\delta}^{{\nu}-1}}{1-{\delta}}


\frac{M}{m}\|x(\cdot,{\alpha})-(A+B)x(\cdot,{\alpha})-g


\|_{\bold C[0,T]},


$$


следовательно, выполняется (15).


\end{pf}





\begin{thebibliography}{9}





\bibitem{1}


Азбелев Н.В., Максимов В.П., Рахматуллина Л.Ф. {\em Введение в теорию


функционально-дифференциальных уравнений}. -- М.: Наука, 1991. -- 280~с.


\bibitem{2}


Azbelev N.V., Maksimov V.P., Rakhmatullina L.F. {\em Introduction to


the theory of linear functional differential equations}. -- Atlanta:


World Federation Publ., 1996. -- 176~p.


\bibitem{3}


Башкиров А.И. {\em Устойчивость решений периодических систем


с последействием}: Дис.\,$\dotsc$


канд. физ.-матем. наук. -- Пермь, 1986. -- 102~с.


\bibitem{4}


Максимов В.П. {\em О формуле Коши для


функционально-дифференциальных уравнений} // Дифференц. уравнения.


-- 1977. -- Т.13. -- \No\,4. -- C.\,601--606.


\bibitem{5}


Максимов В.П.


{\em Линейное функционально-дифференциальное уравнение}: Дис. \,$\dotsc$


канд. физ.-матем. наук. -- Казань, 1974. -- 119~с.


\bibitem{6}


Красносельcкий М.А., Вайникко Г.М., Забрейко П.П.,


Рутицкий Я.Б., Стеценко В.Я.


{\em Приближенное решение операторных уравнений}. -- М.: Наука, 1969.


-- 455~с.





\end{thebibliography}





\vskip 2cm





\post{\it Пермский государственный университет}{\it Поступила}


\post{}{24.05.1999}


\label{end}


\end{document}








